\documentclass[11pt,a4paper]{article}

% --- PACOTES BÁSICOS E TIPOGRAFIA ---
\usepackage[utf8]{inputenc}
\usepackage[T1]{fontenc}
\usepackage[portuguese]{babel}
\usepackage{lmodern}
\usepackage{microtype}

% --- GEOMETRIA E ESPAÇAMENTO ---
\usepackage[top=2.5cm, bottom=2.5cm, left=2.3cm, right=2.3cm, headheight=16pt]{geometry}
\usepackage{setspace}
\onehalfspacing
\usepackage{parskip}

% --- CORES E DESIGN SYSTEM ---
\usepackage{xcolor}
\definecolor{agoraDark}{HTML}{1E293B}      % Slate 800
\definecolor{agoraPrimary}{HTML}{0F172A}   % Slate 900
\definecolor{agoraAccent}{HTML}{EA580C}    % Orange 600 (Cursor/Agora Accent)
\definecolor{agoraBlue}{HTML}{2563EB}      % Blue 600
\definecolor{agoraGreen}{HTML}{16A34A}     % Green 600
\definecolor{agoraPurple}{HTML}{7C3AED}    % Purple 600
\definecolor{bgLight}{HTML}{F8FAFC}        % Slate 50
\definecolor{boxBorder}{HTML}{CBD5E1}      % Slate 300
\definecolor{codeBg}{HTML}{0F172A}         % Dark code background

% --- BOXES MODERNOS ---
\usepackage{tcolorbox}
\tcbset{
    boxrule=0.8pt,
    arc=4pt,
    fonttitle=\bfseries\sffamily,
    coltitle=white,
    colback=bgLight,
    colframe=agoraDark
}

\newtcolorbox{infobox}[2][agoraBlue]{
    colframe=#1,
    colback=bgLight,
    coltitle=white,
    title={#2}
}

\newtcolorbox{warningbox}[2][agoraAccent]{
    colframe=#1,
    colback=bgLight,
    coltitle=white,
    title={#2}
}

\newtcolorbox{successbox}[2][agoraGreen]{
    colframe=#1,
    colback=bgLight,
    coltitle=white,
    title={#2}
}

% --- TABELAS E GRÁFICOS ---
\usepackage{tabularx}
\usepackage{booktabs}
\usepackage{array}
\usepackage{tikz}
\usetikzlibrary{shapes.geometric, arrows.meta, positioning, calc, backgrounds, fit, shadows}

% --- CABEÇALHOS E RODAPÉS ---
\usepackage{fancyhdr}
\pagestyle{fancy}
\fancyhf{}
\fancyhead[L]{\small\sffamily\textcolor{gray}{\textbf{Ágora ENEM} | Documento de Arquitetura \& Engenharia}}
\fancyhead[R]{\small\sffamily\textcolor{gray}{Visão Técnica do Sistema}}
\fancyfoot[C]{\sffamily\thepage}
\renewcommand{\headrulewidth}{0.4pt}
\renewcommand{\footrulewidth}{0.4pt}

% --- TÍTULOS E SEÇÕES ---
\usepackage{titlesec}
\titleformat{\section}
  {\normalfont\Large\bfseries\sffamily\color{agoraDark}}{\thesection}{1em}{}[\color{agoraAccent}\titlerule]
\titleformat{\subsection}
  {\normalfont\large\bfseries\sffamily\color{agoraBlue}}{\thesubsection}{1em}{}
\titleformat{\subsubsection}
  {\normalfont\normalsize\bfseries\sffamily\color{agoraDark}}{\thesubsubsection}{1em}{}

% --- CÓDIGO E LISTAGENS ---
\usepackage{listings}
\lstdefinelanguage{JavaScript}{
  keywords={break, case, catch, class, const, continue, debugger, default, delete, do, else, export, extends, finally, for, function, if, import, in, instanceof, new, return, super, switch, this, throw, try, typeof, var, void, while, with, yield, async, await, let, static, null, true, false, undefined},
  morecomment=[l]{//},
  morecomment=[s]{/*}{*/},
  morestring=[b]',
  morestring=[b]",
  morestring=[b]`
}

\lstset{
    backgroundcolor=\color{bgLight},
    basicstyle=\ttfamily\footnotesize\color{agoraDark},
    keywordstyle=\bfseries\color{agoraAccent},
    stringstyle=\color{agoraGreen},
    commentstyle=\color{gray},
    numberstyle=\tiny\color{gray},
    numbers=left,
    stepnumber=1,
    numbersep=8pt,
    showstringspaces=false,
    breaklines=true,
    frame=single,
    rulecolor=\color{boxBorder},
    tabsize=2
}

% --- LINKS E SUMÁRIO ---
\usepackage[pdftex, hyperfootnotes=false, plainpages=false, pdfpagelabels]{hyperref}
\hypersetup{
    colorlinks=true,
    linkcolor=agoraBlue,
    urlcolor=agoraAccent,
    citecolor=agoraGreen,
    pdftitle={Guia de Arquitetura e Engenharia - Agora ENEM},
    pdfauthor={Equipe Agora ENEM / Clara Silveira}
}
\usepackage{enumitem}

\begin{document}
\pagenumbering{gobble}

% ==============================================================================
% CAPA MODERNA E ELEGANTE
% ==============================================================================
\begin{titlepage}
    \begin{tikzpicture}[remember picture, overlay]
        % Faixa decorativa superior
        \fill[agoraDark] (current page.north west) rectangle ([yshift=-4.5cm]current page.north east);
        \fill[agoraAccent] ([yshift=-4.5cm]current page.north west) rectangle ([yshift=-4.7cm]current page.north east);
        
        % Badge no topo
        \node[anchor=north west, text=white, font=\sffamily\bfseries\small] at ([xshift=2.5cm, yshift=-1.5cm]current page.north west) {PROJETO ÁGORA $\cdot$ PLATAFORMA DE CORREÇÃO PEDAGÓGICA IA};
        \node[anchor=north west, text=white, font=\sffamily\Huge\bfseries] at ([xshift=2.5cm, yshift=-2.4cm]current page.north west) {ÁGORA ENEM};
        \node[anchor=north west, text=orange!30, font=\sffamily\large] at ([xshift=2.5cm, yshift=-3.5cm]current page.north west) {Arquitetura de Software, Engenharia de IA e Guia de Funcionamento};
    \end{tikzpicture}

    \vspace{4.8cm}

    \begin{tcolorbox}[colframe=agoraBlue, title={Resumo Executivo do Documento}]
        Este documento foi concebido para explicar detalhadamente \textbf{o que está por trás} da plataforma \textbf{Ágora ENEM}: sua arquitetura de software, o pipeline de inteligência artificial multimodal com visão computacional, a avaliação cruzada entre a \textbf{Matriz Oficial do ENEM} e a \textbf{Matriz Qualitativa SISEDU / SPAECE (Projeto Ágora Escolar)}, o sistema de fallback resiliente de modelos Google Gemini, a gestão de filas e a interface offline-first voltada à produtividade docente.
    \end{tcolorbox}

    \vspace{0.8cm}

    \begin{center}
        \begin{tabularx}{\textwidth}{r X}
            \textbf{Repositório Oficial:} & \texttt{clarasilveira17/AgoraENEM} \\
            \textbf{Objetivo do Projeto:} & Avaliação Automatizada \& Devolutiva Pedagógica (ENEM $\times$ SISEDU) \\
            \textbf{Linguagens / Stacks:} & Node.js (Express), React 19 (Vite, TailwindCSS), Gemini AI SDK \\
            \textbf{Bancos de Dados:} & Supabase (PostgreSQL Cloud) + SQLite Local + Dexie.js (IndexedDB) \\
            \textbf{Modelos de IA:} & Família Google Gemini Flash (3.5 Flash Lite, 3.7 Flash, Flash Latest) \\
            \textbf{Data da Edição:} & Setembro de 2026 \\
        \end{tabularx}
    \end{center}

    \vfill

    \begin{center}
        \begin{tikzpicture}
            \node[draw=agoraDark, rounded corners=6pt, fill=bgLight, inner sep=10pt, text width=0.85\textwidth, align=center] {
                \textbf{\sffamily Visão Geral:} Da foto da folha manuscrita ao boletim oficial com competências C1--C5, desvios e proposta de intervenção em menos de 10 segundos por aluno.
            };
        \end{tikzpicture}
    \end{center}

    \vspace{0.5cm}
\end{titlepage}

% ==============================================================================
% SUMÁRIO
% ==============================================================================
\tableofcontents
\newpage
\pagenumbering{arabic}

% ==============================================================================
% SEÇÃO 1: VISÃO GERAL E MOTIVAÇÃO
% ==============================================================================
\section{Visão Geral e Contexto Pedagógico}

\subsection{O Desafio Escolar}
No modelo tradicional de ensino, a preparação para a Redação do ENEM (Exame Nacional do Ensino Médio) impõe um desafio colossal aos professores de Língua Portuguesa e Redação:
\begin{itemize}[leftmargin=1.5em]
    \item \textbf{Sobrecarga de Correção:} Corrigir manualmente 200 a 400 redações manuscritas por mês demanda dezenas de horas extraclasse.
    \item \textbf{Latência de Feedback:} Os estudantes frequentemente recebem a devolutiva semanas após a escrita, perdendo o momento pedagógico ideal para assimilar seus erros.
    \item \textbf{Subjetividade e Inconsistência:} Manter a mesma régua de pontuação (critérios das 5 Competências do ENEM de 0 a 200 pontos cada) em pilhas de folhas físicas é exaustivo e sujeito a variações.
\end{itemize}

\subsection{A Solução Ágora ENEM}
A plataforma \textbf{Ágora ENEM} nasceu para solucionar essa dor integrando \textbf{Visão Computacional Multimodal (IA)} e uma \textbf{Interface Human-in-the-Loop}:
\begin{enumerate}[leftmargin=1.5em]
    \item O professor fotografa as folhas oficiais de redação ou o estudante digita seu texto.
    \item A IA transcreve a caligrafia manuscrita integralmente e aplica os critérios do ENEM (C1 a C5) e da instituição (Sisedu).
    \item A plataforma gera uma nota de 0 a 1000 com justificativas por competência, pontos fortes, desvios gramaticais pontuais e o detalhamento da Proposta de Intervenção (Agente, Ação, Meio/Modo, Efeito e Detalhamento).
    \item O professor utiliza a tela de \textbf{Validação Rápida} para auditar a transcrição lado a lado com a foto original e ajustar notas com um único clique antes de liberar o boletim para o aluno.
\end{enumerate}

\begin{successbox}{Human-in-the-Loop (O Professor no Centro)}
A inteligência artificial não substitui o professor: ela atua como um \textbf{super-assistente} que digita, transcreve e pré-avalia com precisão matemática, cabendo ao docente a validação final e o direcionamento humanizado.
\end{successbox}

% ==============================================================================
% SEÇÃO 2: ARQUITETURA GERAL DO SISTEMA
% ==============================================================================
\section{Arquitetura Geral do Sistema}

O Ágora ENEM adota uma arquitetura em camadas orientada a serviços (Service-Oriented Architecture), com separação clara de responsabilidades entre Frontend, Gateway/Backend, Fila de Processamento, Motor de IA e Camada de Persistência Híbrida.

\begin{figure}[htbp]
\centering
\begin{tikzpicture}[
    node distance=1.2cm and 1.8cm,
    font=\sffamily\small,
    box/.style={rectangle, draw=agoraDark, rounded corners=4pt, fill=bgLight, text width=3.2cm, minimum height=1.4cm, align=center, thick},
    aiBox/.style={rectangle, draw=agoraAccent, rounded corners=4pt, fill=orange!5, text width=3.4cm, minimum height=1.5cm, align=center, thick},
    dbBox/.style={rectangle, draw=agoraBlue, rounded corners=4pt, fill=blue!5, text width=3.2cm, minimum height=1.4cm, align=center, thick},
    arrow/.style={-{Stealth[length=3mm]}, thick, draw=agoraDark}
]

    % Frontend
    \node[box, fill=purple!5, draw=agoraPurple] (front) {\textbf{Frontend SPA}\\\scriptsize React 19 + Vite\\Dexie.js (Offline First)};
    
    % Backend Gateway
    \node[box, right=of front] (back) {\textbf{Backend Express}\\\scriptsize Rate Limiting + Auth\\Security Headers};
    
    % Queue
    \node[box, below=of back] (queue) {\textbf{Queue Service}\\\scriptsize Controle 4.5s (13 RPM)\\Retry Backoff};
    
    % AI Engine
    \node[aiBox, right=of back] (ai) {\textbf{AI Service (Gemini)}\\\scriptsize Agente Unificado\\OCR + ENEM C1-C5\\Fallback em Cascata};
    
    % Databases
    \node[dbBox, below=of front] (sqlite) {\textbf{SQLite Local}\\\scriptsize Cache local \& Backup\\Tolerância a falhas};
    \node[dbBox, below=of ai] (supa) {\textbf{Supabase (Cloud)}\\\scriptsize PostgreSQL + JSONB\\Autenticação \& Sync};

    % Conexões
    \draw[arrow] (front) -- node[above]{\scriptsize HTTP/JSON} (back);
    \draw[arrow] (back) -- node[right]{\scriptsize Enfileira} (queue);
    \draw[arrow] (queue) -- node[above]{\scriptsize Executa (Safe RPM)} (ai);
    \draw[arrow] (ai) -- node[above]{\scriptsize Retorna JSON} (back);
    \draw[arrow] (back) -- node[above]{\scriptsize Persiste} (supa);
    \draw[arrow] (back) -- node[left]{\scriptsize Fallback} (sqlite);
    \draw[arrow, dashed] (front) -- node[left]{\scriptsize IndexedDB} (sqlite);

\end{tikzpicture}
\caption{Diagrama de Blocos da Arquitetura Ágora ENEM.}
\label{fig:arquitetura}
\end{figure}

\subsection{Visão das Camadas}
\begin{itemize}[leftmargin=1.5em]
    \item \textbf{1. Camada de Apresentação (Frontend):} Construída em React 19 com compilação ultra-rápida via Vite. Implementa armazenamento local em IndexedDB (Dexie.js), permitindo que o professor trabalhe em escolas sem internet estável.
    \item \textbf{2. Camada de Controle e Segurança (Backend API):} Servidor Node.js (ES Modules) com Express, autenticação JWT, proteção RBAC (\texttt{ADMIN}, \texttt{PROFESSOR}, \texttt{ESTUDANTE}), Correlation IDs (\texttt{X-Request-Id}) e sanitização rigorosa.
    \item \textbf{3. Camada de Fila e Resiliência (\texttt{queueService}):} Gerenciador de concorrência com espaçamento temporal (4.5s) que garante o cumprimento estrito dos limites de taxa gratuitos do Gemini (15 RPM), além de sistema de reprocessamento automático.
    \item \textbf{4. Camada de Inteligência Artificial (\texttt{aiService}):} Motor multimodal que ingere a imagem em Base64, executa a transcrição OCR e a avaliação analítica em um \textbf{Agente Único Unificado}, com fallback automático entre 6 variantes de modelos Gemini.
    \item \textbf{5. Camada de Persistência Híbrida:} Supabase PostgreSQL como banco relacional em nuvem e SQLite local (\texttt{better-sqlite3}) como armazenamento redundante e fallback imediato.
\end{itemize}

% ==============================================================================
% SEÇÃO 3: O MOTOR DE INTELIGÊNCIA ARTIFICIAL
% ==============================================================================
\section{O Motor de Inteligência Artificial Multimodal}

\subsection{Agente Único Unificado: A Grande Otimização}
Nas primeiras versões do protótipo, o sistema realizava duas chamadas consecutivas à API de IA:
\begin{enumerate}
    \item Uma chamada para detecção de cabeçalho e transcrição OCR da caligrafia;
    \item Uma segunda chamada enviando o texto transcrito para avaliação das competências.
\end{enumerate}

\textbf{O Gargalo Identificado:} Esse modelo duplicava a latência, consumia o dobro de tokens e esgotava rapidamente a cota de 15 requisições por minuto (RPM) da API gratuita.

\textbf{A Inovação Implementada:} A plataforma migrou para o \textbf{Agente Único Multimodal} (\texttt{agenteAvaliadorUnificado}). Em uma \textbf{única requisição}, a IA recebe a imagem da folha oficial e realiza simultaneamente:
\begin{itemize}[leftmargin=1.5em]
    \item A leitura ótica dos caracteres (OCR) com transcrição integral e fiel para a chave \texttt{texto\_transcrito};
    \item A extração do Nome do Aluno e Turma no cabeçalho;
    \item A correção pedagógica completa das 5 Competências do ENEM (C1 a C5);
    \item A identificação dos 5 elementos obrigatórios da Proposta de Intervenção (C5);
    \item A pontuação e critérios institucionais complementares (Sisedu).
\end{itemize}

\subsection{Sistema de Fallback em Cascata de Modelos}
Para blindar o sistema contra indisponibilidades de modelos específicos, quotas temporárias ou variações na API do Google, foi criada a função \texttt{generateContentWithFallback}, que itera sequencialmente pela seguinte lista de prioridade:

\begin{lstlisting}[language=JavaScript, caption={Configuração da Cascata de Modelos Gemini}]
export const GEMINI_MODELS = [
  'gemini-3.5-flash-lite',  // 1a Opcao: Alta cota (15 RPM / 500 RPD)
  'gemini-3.7-flash',       // 2a Opcao: Maior capacidade de raciocinio
  'gemini-flash-latest',    // 3a Opcao: Alias dinamico para o ultimo Flash
  'gemini-3.5-flash',       // 4a Opcao: Versao estavel padrao
  'gemini-3.8-flash',       // 5a Opcao: Modelo experimental
  'gemini-3.6-flash'        // 6a Opcao: Fallback final de seguranca
];
\end{lstlisting}

Se um modelo retornar erro \texttt{404 (Not Found)}, \texttt{429 (Rate Limit)} ou \texttt{503 (Service Unavailable)}, o backend captura a exceção imediatamente e reenvia a requisição para o próximo modelo da lista sem quebrar a experiência do usuário.

\subsection{Parser JSON Ultra-Resiliente (\texttt{extractBalancedJson})}
Um dos desafios críticos ao lidar com LLMs (Large Language Models) que retornam JSON estruturado é a ocorrência ocasional de textos explicativos, quebras ou comentários após a chave de fechamento \texttt{\}} do objeto principal.

Para garantir 100\% de taxa de sucesso no parse, o Ágora ENEM implementa um algoritmo de máquina de estados (\texttt{extractBalancedJson}) que:
\begin{itemize}[leftmargin=1.5em]
    \item Rastreia o nível de profundidade de abertura e fechamento de delimitadores (\texttt{\{} e \texttt{\}});
    \item Ignora caracteres dentro de strings e sequências de escape (\texttt{\textbackslash"});
    \item Interrompe a leitura assim que a profundidade retorna a zero, descartando qualquer caractere espúrio subsequente;
    \item Repara automaticamente strings com aspas duplas internas não-escapadas e remove caracteres de controle invisíveis (\texttt{\textbackslash u0000} a \texttt{\textbackslash u001F}).
\end{itemize}

\begin{infobox}{Validação por Testes Unitários}
Essa camada de sanitização e parse possui uma suíte dedicada de 30 testes unitários automatizados no \textbf{Vitest} (\texttt{aiService.test.js}), validando cenários extremos de truncamento, markdown malformado e caracteres residuais.
\end{infobox}

% ==============================================================================
% SEÇÃO 4: A MATRIZ SISEDU E AVALIAÇÃO PEDAGÓGICA QUALITATIVA
% ==============================================================================
\section{A Matriz SISEDU e a Avaliação Pedagógica Qualitativa}

\subsection{O que é o SISEDU no Projeto Ágora Escolar}
Enquanto a matriz oficial do ENEM fornece uma nota somativa de 0 a 1000 pontos voltada a ranqueamento e processos seletivos, o \textbf{SISEDU (Sistema de Avaliação Educacional Integrado)} atua como uma \textbf{matriz diagnóstica e formativa}.

O objetivo do SISEDU no \textbf{Projeto Ágora} é avaliar a maturidade intelectual, a postura cidadã e a capacidade reflexiva do estudante em duas grandes dimensões: \textbf{Discursiva} e \textbf{Ético-Crítica}.

\begin{figure}[htbp]
\centering
\begin{tikzpicture}[
    font=\sffamily\small,
    dimBox/.style={rectangle, draw=agoraDark, rounded corners=4pt, fill=bgLight, text width=6.2cm, minimum height=3.5cm, align=left, thick},
    titleBox/.style={rectangle, fill=agoraDark, text=white, font=\sffamily\bfseries\footnotesize, rounded corners=2pt, inner sep=4pt}
]

    % Caixa Dimensão Discursiva
    \node[dimBox, fill=blue!3, draw=agoraBlue] (disc) at (0,0) {
        \vspace{0.4cm}
        \begin{itemize}[leftmargin=1.2em, itemsep=2pt]
            \item \textbf{Clareza da Tese (D05/D06):} Postura crítica e autoria.
            \item \textbf{Argumentação (D14/D16):} Sustentação lógica e projeto.
            \item \textbf{Repertório (D13):} Conexão interdisciplinar e pertinência.
        \end{itemize}
        \scriptsize\textbf{Níveis:} Inicial $\rightarrow$ Adequado $\rightarrow$ Avançado
    };
    \node[titleBox, fill=agoraBlue, anchor=north west] at (disc.north west) {1. Dimensão Discursiva};

    % Caixa Dimensão Ético-Crítica
    \node[dimBox, fill=orange!3, draw=agoraAccent] (etic) at (7.2,0) {
        \vspace{0.4cm}
        \begin{itemize}[leftmargin=1.2em, itemsep=2pt]
            \item \textbf{Empatia \& Alteridade:} Reconhecimento do outro e direitos.
            \item \textbf{Justificação Moral:} Fundamentação ética e racional.
            \item \textbf{Conclusão Propositiva:} Viabilidade transformadora.
        \end{itemize}
        \scriptsize\textbf{Níveis:} Inicial $\rightarrow$ Adequado $\rightarrow$ Avançado
    };
    \node[titleBox, fill=agoraAccent, anchor=north west] at (etic.north west) {2. Dimensão Ético-Crítica};

\end{tikzpicture}
\caption{Estrutura das Duas Dimensões Avaliativas da Matriz SISEDU.}
\label{fig:sisedu_matriz}
\end{figure}

\subsection{Rubricas das Dimensões SISEDU}

\subsubsection{1. Dimensão Discursiva}
Avalia como o estudante estrutura o pensamento e constrói seu posicionamento:
\begin{itemize}[leftmargin=1.5em]
    \item \textbf{Tese e Autoria:}
    \begin{itemize}
        \item \textit{Inicial:} Opinião vaga, implícita ou restrita à repetição da proposta.
        \item \textit{Em Desenvolvimento / Adequado:} Ponto de vista explícito e delimitado na introdução.
        \item \textit{Avançado:} Tese problematizadora com forte marca de autoria e posicionamento crítico maduro.
    \end{itemize}
    \item \textbf{Progressão Argumentativa:}
    \begin{itemize}
        \item \textit{Inicial:} Exemplos isolados, senso comum e fragmentação temática.
        \item \textit{Em Desenvolvimento / Adequado:} Relações de causa, consequência e projeto de texto coerente.
        \item \textit{Avançado:} Argumentação sólida, contra-argumentação eficiente e articulação interdisciplinar.
    \end{itemize}
    \item \textbf{Uso do Repertório Sociocultural:}
    \begin{itemize}
        \item \textit{Inicial:} Ausência de referências ou citações baseadas exclusivamente nos textos motivadores.
        \item \textit{Em Desenvolvimento / Adequado:} Menção a fatos históricos ou leis de forma pertinente.
        \item \textit{Avançado:} Repertório produtivo integrado organicamente à argumentação (filosofia, sociologia, arte e dados).
    \end{itemize}
\end{itemize}

\subsubsection{2. Dimensão Ético-Crítica}
Avalia o compromisso social e o alcance ético do estudante diante da realidade brasileira:
\begin{itemize}[leftmargin=1.5em]
    \item \textbf{Empatia e Alteridade (Direitos Humanos):} Capacidade de compreender a vulnerabilidade de grupos marginalizados sem posturas punitivas ou preconceituosas.
    \item \textbf{Justificação Moral / Axiológica:} Fundamentação de escolhas argumentativas em princípios de justiça social, equidade e solidariedade democrática.
    \item \textbf{Eficácia Propositiva (Conclusão Crítica):} Capacidade de articular ações viáveis com poder de transformação social real.
\end{itemize}

\subsection{Integração com a Matriz de Descritores SPAECE (Ceará)}
A plataforma correlaciona a avaliação discursiva aos descritores de proficiência do **SPAECE** (Sistema Permanente de Avaliação da Educação Básica do Ceará), permitindo diagnósticos institucionais precisos:

\begin{center}
\begin{tabularx}{\textwidth}{l l p{8.5cm}}
\toprule
\textbf{Código} & \textbf{Competência} & \textbf{Descrição Pedagógica do Descritor} \\
\midrule
\texttt{D05} & Leitura / Tese & Identificar o tema central e a tese defendida no texto. \\
\texttt{D06} & Argumentação & Distinguir um fato da opinião ou juízo de valor a ele associado. \\
\texttt{D12} & Coesão Lexical & Estabelecer relações lógico-discursivas marcadas por conectores e operadores. \\
\texttt{D13} & Recursos Expressivos & Avaliar os efeitos de sentido decorrentes da pontuação e figuras de linguagem. \\
\texttt{D14} & Hierarquia Textual & Diferenciar informação principal de argumentos secundários de apoio. \\
\texttt{D15} & Continuidade Textual & Reconhecer mecanismos de retomada anafórica, elipses e contra-argumentação. \\
\texttt{D16} & Intertextualidade & Identificar o diálogo intertextual e posicionamentos contrastantes. \\
\bottomrule
\end{tabularx}
\end{center}

\subsection{Plano de Intervenção Pedagógica Automatizado}
Sempre que a IA detecta que o aluno atingiu o nível \textbf{Inicial} em qualquer dos descritores SISEDU ou competências do ENEM, o sistema injeta no boletim um **Plano de Intervenção Pedagógica Personalizado**, por exemplo:
\begin{quote}
\small\sffamily\textit{``Atenção: O estudante apresentou descritores em Nível Inicial. Recomenda-se aplicar atividade direcionada de reescrita com foco em contra-argumentação (D15) e ampliação de repertório sociocultural legítimo (D13).''}
\end{quote}

O professor pode optar por imprimir o boletim apenas com as notas do ENEM, apenas com a grade qualitativa SISEDU, ou com a \textbf{visão integrada completa}.

% ==============================================================================
% SEÇÃO 5: ENGENHARIA DE BACKEND E FILA
% ==============================================================================
\section{Engenharia de Backend \& Fila de Processamento}

\subsection{Estrutura de Diretórios e Modularização}
O backend é estruturado sob o padrão arquitetural em camadas:

\begin{center}
\begin{tabularx}{\textwidth}{l p{10cm}}
\toprule
\textbf{Diretório / Arquivo} & \textbf{Responsabilidade Técnica} \\
\midrule
\texttt{server.js} & Ponto de entrada, middlewares globais, segurança e graceful shutdown. \\
\texttt{src/config/} & Configuração de clientes: Supabase, Gemini SDK e SQLite. \\
\texttt{src/controllers/} & Orquestração HTTP, validação de payload e respostas da API. \\
\texttt{src/services/} & Regras de negócio: Fila (\texttt{queueService}) e IA (\texttt{aiService}). \\
\texttt{src/repositories/} & Camada de acesso a dados (CRUD de redações e usuários). \\
\texttt{src/middleware/} & Autenticação JWT, controle de acesso RBAC e Rate Limiters. \\
\texttt{src/utils/} & Logger contextualizado, cálculo de turmas e sanitização. \\
\bottomrule
\end{tabularx}
\end{center}

\subsection{Controle de Fluxo e Fila (\texttt{RateLimitedQueueService})}
O envio em lote (por exemplo, 40 fotos de uma turma de uma só vez) excederia instantaneamente o teto de 15 RPM da API Gemini caso todas fossem disparadas em paralelo.

Para solucionar isso, o \texttt{RateLimitedQueueService} implementa:
\begin{enumerate}[leftmargin=1.5em]
    \item \textbf{Espaçamento Temporal Estrito (4500ms):} Entre cada requisição à IA, a fila aguarda 4,5 segundos ($\approx 13$ requisições por minuto), mantendo a taxa sempre abaixo do limite de 15 RPM.
    \item \textbf{Mecanismo de Retry com Backoff Exponencial:} Em caso de falha transitória de rede, o item é re-enfileirado com atraso crescente (10s na primeira tentativa, 20s na segunda).
    \item \textbf{Status \texttt{ERRO\_PROCESSAMENTO} e Worker de Resgate:} Se o item esgotar as 3 tentativas, a imagem e os metadados são salvos no banco com status de erro. Um worker em background roda a cada 3 minutos para reprocessar automaticamente itens pendentes.
\end{enumerate}

\begin{lstlisting}[language=JavaScript, caption={Loop de Processamento com Rate Limiter da Fila}]
// Trecho de execucao com espacamento temporal seguro
while (this.queue.length > 0) {
  const current = this.queue.shift();
  current.status = 'PROCESSANDO';
  
  try {
    const result = await this.executeEvaluation(current.item, current.requestingUser);
    current.status = 'CONCLUIDO';
  } catch (error) {
    this.handleFailure(current, error);
  }
  
  // Atraso obrigatorio para blindagem contra 429 Too Many Requests
  await new Promise(resolve => setTimeout(resolve, this.rateLimitDelayMs));
}
\end{lstlisting}

% ==============================================================================
% SEÇÃO 6: BANCO DE DADOS E MODELO DE DADOS
% ==============================================================================
\section{Esquema do Banco de Dados \& Persistência Híbrida}

A plataforma utiliza o banco de dados relacional **Supabase (PostgreSQL)** com suporte avançado a dados semiestruturados via colunas **JSONB**.

\subsection{Esquema Relacional (PostgreSQL / Supabase)}

\begin{lstlisting}[language=SQL, caption={Estrutura das Tabelas Principais}]
-- 1. Tabela de Usuarios (Controle de Acesso RBAC)
CREATE TABLE IF NOT EXISTS public.users (
  id BIGSERIAL PRIMARY KEY,
  nome TEXT NOT NULL,
  email TEXT UNIQUE NOT NULL,
  senha_hash TEXT NOT NULL,
  role TEXT NOT NULL DEFAULT 'ESTUDANTE', -- 'ADMIN', 'PROFESSOR', 'ESTUDANTE'
  turma TEXT,
  created_at TIMESTAMPTZ DEFAULT NOW()
);

-- 2. Tabela de Redacoes
CREATE TABLE IF NOT EXISTS public.redacoes (
  id BIGSERIAL PRIMARY KEY,
  user_id BIGINT REFERENCES public.users(id) ON DELETE SET NULL,
  nome_aluno TEXT,
  turma_aluno TEXT,
  nome_detectado INTEGER DEFAULT 0,
  data_captura TIMESTAMPTZ DEFAULT NOW(),
  tipo_input TEXT DEFAULT 'imagem',       -- 'imagem' ou 'texto'
  imagem_base64 TEXT,                     -- Armazena foto da folha oficial
  texto_digitado TEXT,                    -- Texto digitado ou transcrito
  is_synced INTEGER DEFAULT 1,
  extracted_data JSONB DEFAULT '{}'::jsonb, -- JSON integral da avaliacao da IA
  nota_final INTEGER,                     -- Pontuacao 0 a 1000
  status_validacao TEXT NOT NULL DEFAULT 'VALIDADA', -- 'VALIDADA', 'PENDENTE'
  validado_por BIGINT REFERENCES public.users(id) ON DELETE SET NULL,
  data_validacao TIMESTAMPTZ
);
\end{lstlisting}

\subsection{Estrutura do Campo \texttt{extracted\_data} (JSONB)}
O campo \texttt{extracted\_data} contém a riqueza analítica gerada pela IA:
\begin{itemize}[leftmargin=1.5em]
    \item \texttt{texto\_transcrito}: Transcrição completa com quebra de parágrafos.
    \item \texttt{competencias}: Objeto contendo C1, C2, C3, C4 e C5 (cada qual com sua pontuação de 0 a 200 e justificativa pedagógica textual).
    \item \texttt{proposta\_intervencao}: Breakdown dos 5 componentes da C5 (Agente, Ação, Meio/Modo, Efeito e Detalhamento).
    \item \texttt{desvios}: Lista de inadequações gramaticais com trecho citado, regra violada e sugestão de reescrita.
    \item \texttt{pontos\_fortes}: Elogios às passagens bem argumentadas e repertórios legítimos.
    \item \texttt{analise\_sisedu}: Critérios da instituição (adequação ao tema, tipologia textual).
\end{itemize}

% ==============================================================================
% SEÇÃO 7: FRONTEND E EXPERIÊNCIA DO USUÁRIO
% ==============================================================================
\section{Engenharia de Frontend \& Telas do Sistema}

\subsection{Design System e Filosofia de Interface}
A interface gráfica foi construída sob a inspiração do ecossistema editorial moderno (estilo \textbf{Cursor}):
\begin{itemize}[leftmargin=1.5em]
    \item \textbf{Cores Harmônicas:} Fundo neutro suave (\texttt{\#f7f7f4}), tipografia nítida (\texttt{\#26251e}) e destaques em Laranja Accent (\texttt{\#f54e00}), Pêssego e Menta.
    \item \textbf{Fontes Especializadas:} \textit{Outfit} e \textit{Inter} para interface e navegação, e \textit{JetBrains Mono} com numeração de linhas para visualização da redação transcrita.
\end{itemize}

\subsection{Mapeamento das Principais Vistas (Views)}
\begin{center}
\begin{tabularx}{\textwidth}{l p{10.5cm}}
\toprule
\textbf{Componente React} & \textbf{Funcionalidade e Experiência do Usuário} \\
\midrule
\texttt{DashboardView} & Painel analítico executivo com KPIs, gráfico de radar das 5 competências do ENEM, média por turma e histórico recente. \\
\texttt{UploaderView} & Área de ingestão de redações: suporta arrastar-e-soltar de dezenas de fotos simultâneas, captura via câmera ou digitação direta de texto. \\
\texttt{ValidacaoRapidaView} & Tela em tela dividida (split-screen): à esquerda a foto original com zoom/pan, à direita o texto transcrito e sliders de ajuste de notas para auditoria em alta velocidade. \\
\texttt{GestaoRedacoesView} & Repositório central de dados com busca em tempo real, filtros por turma, status de validação, faixa de notas e exportação para CSV/Excel. \\
\texttt{ModalDetalhesRedacao} & Modal pedagógico completo com detalhamento dos 5 elementos da intervenção, lista de desvios e comparativo de nota. \\
\texttt{FolhaOficialRedacao} & Visualização que replica fielmente a folha de redação padrão do INEP/ENEM com 30 linhas numeradas. \\
\texttt{DiagnosticoFeiraView} & Relatório de inteligência pedagógica para a coordenação, identificando as principais fragilidades gramaticais e temáticas da escola. \\
\texttt{RankingView} & Mural de desempenho dos estudantes com medalhas e classificação por turma. \\
\bottomrule
\end{tabularx}
\end{center}

% ==============================================================================
% SEÇÃO 8: FLUXO DE VIDA DE UMA REDAÇÃO
% ==============================================================================
\section{Ciclo de Vida de uma Redação (Fluxo Ponta a Ponta)}

Abaixo detalha-se o ciclo completo percorrido por uma redação desde a folha física até o relatório final:

\begin{figure}[htbp]
\centering
\begin{tikzpicture}[
    font=\sffamily\footnotesize,
    stepBox/.style={rectangle, draw=agoraDark, rounded corners=3pt, fill=bgLight, text width=2.8cm, minimum height=1.3cm, align=center, thick},
    arrow/.style={-{Stealth[length=2.5mm]}, thick, draw=agoraAccent}
]

    \node[stepBox] (s1) {\textbf{1. Captura}\\\scriptsize Foto ou Texto digitado pelo Professor/Aluno};
    \node[stepBox, right=0.8cm of s1] (s2) {\textbf{2. Offline Store}\\\scriptsize Salva no Dexie.js (IndexedDB local)};
    \node[stepBox, right=0.8cm of s2] (s3) {\textbf{3. Enfileiramento}\\\scriptsize Envio ao backend e fila com rate limit};
    
    \node[stepBox, below=1.2cm of s3, fill=orange!5, draw=agoraAccent] (s4) {\textbf{4. IA Multimodal}\\\scriptsize OCR + C1-C5 + Intervencao no Gemini};
    \node[stepBox, left=0.8cm of s4, fill=blue!5, draw=agoraBlue] (s5) {\textbf{5. Persistencia}\\\scriptsize Gravacao no Supabase e SQLite};
    \node[stepBox, left=0.8cm of s5, fill=green!5, draw=agoraGreen] (s6) {\textbf{6. Validacao \& PDF}\\\scriptsize Auditoria docente e emissao do Boletim};

    \draw[arrow] (s1) -- (s2);
    \draw[arrow] (s2) -- (s3);
    \draw[arrow] (s3) -- (s4);
    \draw[arrow] (s4) -- (s5);
    \draw[arrow] (s5) -- (s6);

\end{tikzpicture}
\caption{Fluxograma do Ciclo de Vida de Processamento da Redação.}
\label{fig:ciclo_vida}
\end{figure}

% ==============================================================================
% SEÇÃO 9: GUIA PRÁTICO DE EXECUÇÃO
% ==============================================================================
\section{Guia de Inicialização e Operação do Sistema}

\subsection{Execução Simplificada no Windows}
O projeto conta com scripts utilitários na raiz:
\begin{itemize}[leftmargin=1.5em]
    \item \texttt{iniciar.bat} ou \texttt{iniciar.exe}: Inicializa simultaneamente o servidor backend na porta \texttt{3001} e a aplicação frontend Vite na porta \texttt{5173}.
\end{itemize}

\subsection{Execução Manual via Terminal}
Para inicializar os serviços individualmente:

\begin{lstlisting}[language=bash, caption={Comandos para Inicializacao}]
# 1. No diretorio /backend
cd backend
npm install
npm run dev      # Inicia o servidor Express na porta 3001

# 2. No diretorio /frontend (em outro terminal)
cd frontend
npm install
npm run dev      # Inicia o servidor Vite na porta 5173
\end{lstlisting}

\subsection{Variáveis de Ambiente (\texttt{backend/.env})}
\begin{lstlisting}[language=bash, caption={Variaveis de Ambiente Obrigatorias}]
PORT=3001
GEMINI_API_KEY=AIzaSy...              # Chave de API do Google AI Studio
SUPABASE_URL=https://xyz.supabase.co  # URL do projeto Supabase
SUPABASE_KEY=eyJhbGciOi...            # Service Role / Anon Key
JWT_SECRET=chave_secreta_jwt          # Assinatura segura de tokens
\end{lstlisting}

% ==============================================================================
% SEÇÃO 10: CONCLUSÃO E ROADMAP
% ==============================================================================
\section{Conclusão e Próximos Passos}

O projeto \textbf{Ágora ENEM} estabelece um novo padrão na avaliação educacional orientada por Inteligência Artificial ao conciliar:
\begin{itemize}[leftmargin=1.5em]
    \item \textbf{Eficiência Extrema:} Redução de 90\% no tempo de correção de redações manuscritas.
    \item \textbf{Resiliência e Tolerância a Falhas:} Fallback de modelos de IA, parser balanceado de JSON e fila com espaçamento seguro de RPM.
    \item \textbf{Rigor Pedagógico:} Análise integral segundo as 5 Competências do ENEM e matriz do Sisedu.
    \item \textbf{Soberania Docente:} O professor permanece com a palavra final através de ferramentas visuais ágeis de auditoria.
\end{itemize}

\vspace{1cm}
\begin{center}
    \rule{0.6\textwidth}{0.4pt} \\
    \vspace{0.3cm}
    \small\sffamily\textcolor{gray}{\textbf{Ágora ENEM} --- Tecnologia \& Educação a Serviço do Futuro.}
\end{center}

\end{document}
